\section{Kausale Überprüfung mit Activation Patching}

Alle bisherigen Verfahren liefern Korrelationen. Sie zeigen, dass eine Information vorhanden ist oder genutzt zu werden scheint, aber nicht, dass eine bestimmte Komponente für die Vorhersage tatsächlich notwendig ist. Diesen Nachweis liefert das \textbf{Activation Patching}, das in der mechanistischen Interpretierbarkeit als kausaler Goldstandard gilt \cite{meng2022locating, goldowskydill2023localizing}. Das Prinzip ist ein kontrolliertes Experiment, bei dem ein \emph{sauberer} und ein \emph{korrumpierter} Durchlauf verwendet werden, die sich nur im sentimenttragenden Wort unterscheiden. Im korrumpierten Durchlauf wird dann gezielt die \gls{activation} einer einzelnen Schicht an einer einzelnen Token-Position durch die zuvor gespeicherte saubere Version ersetzt. Verändert sich die Vorhersage dadurch wieder in Richtung des sauberen Ergebnisses, enthält genau diese Schicht und Position die entscheidende Information. Normalisiert bedeutet ein Effekt von 0 keine Wiederherstellung, ein Effekt von 1 eine vollständige Wiederherstellung der sauberen Vorhersage. Technisch läuft das über PyTorch Forward Hooks. Zunächst werden alle sauberen \glspl{activation} zwischengespeichert, danach wird systematisch über jede Kombination aus Schicht und Position gepatcht und der jeweilige Effekt gemessen, und zuletzt werden die Ergebnisse als Heatmap dargestellt und die einflussreichsten Punkte herausgefiltert.

Als Prompt-Paar dient erneut das bereits aus den vorigen Kapiteln bekannte Beispiel \texttt{The food was delicious/disgusting and I felt very} mit dem Zielwort „satisfied“. Dieselbe Formulierung wurde bereits im Rahmen der Logit-Lens- und Attention-Analyse untersucht, sodass sich die Ergebnisse der drei Methoden direkt miteinander vergleichen lassen. Im sauberen Durchlauf erreicht das Zielwort einen Logit von 15,93 bei einer Wahrscheinlichkeit von 10,8~Prozent. Im korrumpierten Durchlauf sinkt der Logit auf 9,54, während die Wahrscheinlichkeit lediglich 0,018~Prozent beträgt. Die deutliche Differenz zeigt, dass die Vorhersage stark vom sentimenttragenden Wort beeinflusst wird. Auf Grundlage der Ergebnisse der vorigen Kapitel war zu erwarten, dass die stärksten kausalen Effekte in den oberen Schichten, insbesondere zwischen Schicht~17 und~23, auftreten würden, da dort die deutlichsten Muster in der Logit-Lens- und Attention-Analyse sichtbar waren. Die nachfolgende Activation-Patching-Analyse zeigt jedoch ein differenzierteres Bild und legt nahe, dass die entscheidende Sentimentinformation bereits deutlich früher im Modell kodiert wird.

\begin{figure}[htbp]
  \centering
  \includegraphics[width=\textwidth]{activation_patching_heatmap.png}
  \caption{Normalisierter Activation-Patching-Effekt über alle Schichten und Token-Positionen.}
  \label{fig:patch}
\end{figure}
Eine genauere Betrachtung der Heatmap zeigt zudem eine klare Verschiebung der relevanten Information über die Schichten hinweg. In den frühen Schichten von etwa Schicht~0 bis~10 konzentrieren sich die stärksten Patching-Effekte nahezu ausschließlich auf das sentimenttragende Wort „delicious“ an Position~3. Die Sentimentinformation scheint daher zunächst lokal in der Repräsentation dieses Tokens gespeichert zu sein. In den mittleren Schichten zwischen etwa Schicht~10 und~15 verteilt sich der Effekt zunehmend auf nachfolgende Tokenpositionen. In den oberen Schichten ab ungefähr Schicht~16 treten die stärksten Effekte schließlich vor allem bei den Wörtern „felt“ und „very“ auf. Dies deutet darauf hin, dass die ursprünglich im Sentimentwort kodierte Information schrittweise über die Sequenz weitergegeben wird, bis sie an der Vorhersageposition für das Zielwort „satisfied“ genutzt werden kann.

\begin{table}[htbp]
  \centering
  \caption{Die fünf einflussreichsten Patching-Punkte.}
  \label{tab:patchtop}
  \begin{tabular}{lrrlr}
    \toprule
    Rang & Schicht & Position & Token & Effekt \\
    \midrule
    1 & 2 & 3 & delicious & 1,023 \\
    2 & 1 & 3 & delicious & 1,021 \\
    3 & 0 & 3 & delicious & 1,017 \\
    4 & 3 & 3 & delicious & 1,016 \\
    5 & 6 & 3 & delicious & 1,016 \\
    \bottomrule
  \end{tabular}
\end{table}

Patching-Punkte liegen sämtlich an Position~3, also direkt am sentimenttragenden Wort „delicious“, und treten bereits in den frühen Schichten~0 bis~6 auf (Tabelle~5). Der stärkste Effekt findet sich in Schicht~2. Die entscheidende Sentimentinformation wird damit bereits in den ersten Schichten kodiert, während spätere Schichten diese Information offenbar vor allem weiterverarbeiten und für die finale Vorhersage nutzbar machen.

Ergänzend wird das Activation Patching auf Ebene einzelner Attention-Köpfe durchgeführt. Anstatt die gesamte Aktivierung einer Schicht zu ersetzen, wird dabei ausschließlich der Beitrag eines einzelnen Attention-Kopfes an Position~3, also am sentimenttragenden Wort „delicious“, gepatcht (Abbildung~9). Ziel dieser Analyse ist es zu untersuchen, ob sich die zuvor beobachteten starken Effekte auf Schichtebene auf einen bestimmten Kopf zurückführen lassen.

Die Ergebnisse zeigen jedoch ein deutlich anderes Bild als beim Layer-Level-Patching. Der stärkste Einzelkopf befindet sich in Schicht~5, Kopf~12, und erreicht lediglich einen Effekt von 0,019. Dieser Wert liegt mehr als eine Größenordnung unter den stärksten Effekten auf Schichtebene, die Werte von über~1 erreichen. Selbst der einflussreichste Kopf kann somit nur einen sehr kleinen Teil des zuvor beobachteten Effekts erklären.

Auch die Heatmap liefert keinen Hinweis auf einen einzelnen dominanten Kopf. Statt eines klar abgegrenzten Hotspots zeigen sich lediglich schwache und über mehrere Köpfe verteilte Effekte. Dies spricht dagegen, dass die sentimentrelevante Information in einem einzelnen spezialisierten Attention-Kopf gespeichert ist. Vielmehr scheint sie über mehrere Köpfe und Komponenten innerhalb der jeweiligen Schichten verteilt repräsentiert zu werden.

Besonders auffällig ist zudem, dass der in der Attention-Analyse aus Kapitel~6 hervorgehobene Kopf aus Schicht~6, Kopf~3, hier keinen nennenswerten kausalen Einfluss zeigt. Dies verdeutlicht den Unterschied zwischen Attention-Analyse und Activation Patching: Hohe Attention-Gewichte zeigen zwar, welche Verbindungen das Modell nutzt, sie garantieren jedoch nicht, dass der betreffende Kopf auch die für die Vorhersage entscheidende Information enthält. Die Head-Level-Patching-Analyse liefert daher keinen Hinweis auf einen einzelnen „Sentiment Head“, sondern spricht für eine verteilte Repräsentation der sentimentrelevanten Information innerhalb des Modells.
\begin{figure}[htbp]
  \centering
  \includegraphics[width=0.7\textwidth]{activation_patching_head_heatmap.png}
  \caption{Patching auf Ebene einzelner Attention Heads an Position 3 (Token „delicious“).}
  \label{fig:patchhead}
\end{figure}

